\chapter{Installation and database creation}
\label{ops:installation}

\section{Choose the implementation you intend to run}

MiniSQL is a native database server written in MiniLang. The server owns a
database directory; clients connect to its TCP listener and send SQL. The
command-line client is a separate program, so starting a client does not start
a server automatically. The native Windows Workbench is another client of the
same server. Python and Java applications can connect through the supplied
drivers.

This manual describes the development revision printed on its title page.
That revision uses database, table, index, and row format version 2. Its WAL
record framing remains version 1. The application version string and the
storage-format version answer different questions: a binary identifying
itself as MiniSQL 1.1.1 is not sufficient evidence that it can open a
particular database. Record the source revision and build target alongside
the executable.

\warning{Existing databases}{The current source rejects format-1 databases.
Do not point a newly built format-2 server at the only copy of an older
database. Retain the matching old executable, export the schema and data
through that executable, create a new database with the current executable,
and load the exported information. The page-size migration tool does not
implement a format-1 to format-2 conversion.}

\sourcecode{src/minisql/common/format_version.ml}{Persisted format identifiers.}
\sourcecode{docs/formats/FORMAT_COMPATIBILITY.md}{Compatibility policy for the current tree.}

\section{Platforms, prerequisites, and program inventory}
\index{installation}\index{Windows}\index{Linux}\index{build prerequisites}

The supported native targets are Windows x64 and Linux x64. Linux requires
glibc; native TLS requires OpenSSL 3, including \path{libssl.so.3} and
\path{libcrypto.so.3}. Windows uses Schannel and CryptoAPI for native TLS and
needs an operating-system configuration capable of the required TLS 1.3
profile. The Windows GUI is not a portable Linux application.

For a source build, use MiniLangCompilerPy or MiniLangCompilerML 1.3.0 or
newer, as required by the current repository. The Python compiler and the
repository's build/test tooling require Python 3.11 or newer. The Python
database connector itself has a separate Python 3.10 minimum. PowerShell is
needed for the supplied \path{build.ps1} launcher. A direct compiler command
can be used on Linux without that launcher.

\begin{longtable}{p{.27\linewidth}p{.66\linewidth}}
\toprule
Program & Purpose \\
\midrule
\endhead
\path{minisqld} & Creates databases, administers local passwords and encryption,
and runs the server. \\
\path{minisql} & Interactive SQL shell, single-query client, script runner,
and connection checks. \\
\path{minisql-check} & Opens a stopped database and verifies authoritative
storage, rows, overflow values, and indexes. \\
\path{minisql-backup} & Base backups, encrypted backups, restores, WAL archives,
point-in-time recovery, and standby administration. \\
\path{minisql-migrate} & Checks whether a page-size change is needed and
performs a separate offline copy-and-rewrite when requested. \\
\path{minisql-admin} & Windows-only graphical MiniSQL Workbench. \\
\bottomrule
\end{longtable}

Windows names end in \path{.exe}. Windows builds are placed in
\path{build/bin}; Linux builds are placed in \path{build/bin-linux} and omit
that suffix. Keep all programs from the same source revision together. A
checker or backup program is part of the storage implementation, so mixing
it with an unrelated server build is an unnecessary compatibility risk.

\sourcecode{build.ps1}{Target selection, compiler discovery, include paths, and executable inventory.}

\section{Build from source on Windows}
\index{build.ps1}\index{compiler}

Keep sibling checkouts for MiniSQL and the compiler. In PowerShell, change to
the MiniSQL repository and run:

\begin{lstlisting}
git clone https://github.com/MiniLangProject/MiniSQL.git
git clone https://github.com/MiniLangProject/MiniLangCompilerPy.git
cd MiniSQL
$compiler = "..\MiniLangCompilerPy\mlc_win64.py"
.\build.ps1 -Compiler $compiler -Target windows-x64 -AppsOnly
.\build\bin\minisqld.exe --version
\end{lstlisting}

The compiler path can be absolute. If \code{-Compiler} is omitted, the script
checks \code{MINILANG\_COMPILER} and several conventional repository-relative
locations. Explicit paths make a build easier to reproduce. Use
\code{-Python} to select a particular Python executable when more than one
is installed.

\begin{lstlisting}
.\build.ps1 -Compiler C:\Tools\MiniLangCompilerPy\mlc_win64.py `
  -Python C:\Python311\python.exe -Target windows-x64 -AppsOnly
\end{lstlisting}

With the native self-hosted compiler, supply its executable instead:

\begin{lstlisting}
.\build.ps1 -Compiler ..\MiniLangCompilerML\build\mlc_win64.exe `
  -Target windows-x64 -AppsOnly
\end{lstlisting}

The launcher discovers the compiler's \path{std/} include root and selects
its object pipeline for the large self-hosted build. \code{-AppsOnly} builds
the public programs; omitting it also compiles the native test applications.
\code{-Clean} removes the selected binary output directory before rebuilding.
The script verifies both the compiler's exit code and the existence of each
expected executable. A zero exit status from the compiler without an output
file is treated as failure.

Building Linux targets from the same checkout is possible by selecting
\code{-Target linux-x64}. Compiling a Linux ELF file does not execute its
tests on Windows; native execution requires Linux or an appropriate WSL
environment.

\section{Build and run on Linux}
\index{Linux!building}

If PowerShell is available on Linux, the same launcher can select
\code{linux-x64}. The following direct Python-compiler commands show the
equivalent source/include relationship without depending on PowerShell.
Run them from the MiniSQL repository with a sibling compiler checkout:

\begin{lstlisting}[language=bash]
mkdir -p build/bin-linux
compiler=../MiniLangCompilerPy/mlc_win64.py
compiler_root=../MiniLangCompilerPy
python3 "$compiler" src/apps/minisqld/main.ml \
  build/bin-linux/minisqld -I src -I "$compiler_root" \
  --target linux-x64
python3 "$compiler" src/apps/minisql/main.ml \
  build/bin-linux/minisql -I src -I "$compiler_root" \
  --target linux-x64
python3 "$compiler" src/apps/minisql_check/main.ml \
  build/bin-linux/minisql-check -I src -I "$compiler_root" \
  --target linux-x64
python3 "$compiler" src/apps/minisql_backup/main.ml \
  build/bin-linux/minisql-backup -I src -I "$compiler_root" \
  --target linux-x64
python3 "$compiler" src/apps/minisql_migrate/main.ml \
  build/bin-linux/minisql-migrate -I src -I "$compiler_root" \
  --target linux-x64
chmod +x build/bin-linux/minisql*
build/bin-linux/minisqld --version
\end{lstlisting}

The output remains native machine code: Python is a build dependency here,
not an interpreter that runs each SQL query. The optional monitoring,
replication, and qualification scripts use Python independently of the native
server.

If a generated program reports that a shared library cannot be loaded,
inspect the Linux runtime dependencies before changing SQL or database files.
If the compiler cannot find \code{std} modules, check the second include root.
If the shell reports permission denied for a compiled binary, check its
executable permission and the mount's execution policy.

\section{Create the first database}

\subsection{A database is created by the server tool}

MiniSQL does not implement SQL \code{CREATE DATABASE}, \code{DROP DATABASE},
or \code{USE}. A server instance is opened on one database path. Creating a
table in an existing database and creating that database directory are
different operations.

\begin{lstlisting}
# Windows PowerShell, from the repository root
.\build\bin\minisqld.exe --init .\data learning 4096
\end{lstlisting}

\begin{lstlisting}[language=bash]
# Linux, from the repository root
build/bin-linux/minisqld --init ./data learning 4096
\end{lstlisting}

The tool reports a generated path of the form
\path{data/db_<32-hexadecimal-digits>}, the logical name \code{learning},
and the page size. Copy the actual path from this output. The database is not
necessarily stored in a directory named \path{learning}; the UUID-derived
directory avoids confusing a logical database name with its physical identity.

Supported page sizes are 4096, 8192, 16384, and 32768 bytes. Omitting the last
argument selects 4096. Start with 4096 unless a measured workload gives a
reason to choose differently. Page size is persisted when the database is
created. Editing a JSON default later does not change existing pages.

The public \code{--init} command accepts the data root, logical name, and
optional page size. It uses the compiled default database settings. There is
no \code{--init-config} command that consumes arbitrary creation defaults
from JSON. Do not assume that a configuration file is read automatically.

\sourcecode{src/apps/minisqld/main.ml}{initializeDatabase and the --init argument branch.}
\sourcecode{src/minisql/catalog/catalog.ml}{Database creation and UUID-based paths.}

\subsection{Start a local learning server}

Replace the example path below with the path printed by initialization:

\begin{lstlisting}
$database = ".\data\db_REPLACE_WITH_THE_PRINTED_ID"
.\build\bin\minisqld.exe --serve $database 7432 32
\end{lstlisting}

\begin{lstlisting}[language=bash]
database='./data/db_REPLACE_WITH_THE_PRINTED_ID'
build/bin-linux/minisqld --serve "$database" 7432 32
\end{lstlisting}

The process remains running. Its startup banner reports the mode, physical
database path, address, port, connection limit, and request budget. This
particular mode binds to \code{127.0.0.1} and trusts local clients. Every
local client that can reach the port is administrative; a database password
does not protect this mode.

Open a second terminal in the repository and connect:

\begin{lstlisting}
.\build\bin\minisql.exe --ping 7432
.\build\bin\minisql.exe --shell 7432
\end{lstlisting}

On Linux the equivalents are \path{build/bin-linux/minisql --ping 7432} and
\path{build/bin-linux/minisql --shell 7432}. At the SQL prompt:

\begin{lstlisting}[language=SQL]
SELECT 2 + 3 AS answer;
SHOW TABLES;
\end{lstlisting}

The first query should return one row containing 5. A fresh database has no
user-created tables. The client session and the server process have separate
lifetimes: leaving the shell disconnects that client but keeps the server
running. To request a graceful stop from an authorized session, execute:

\begin{lstlisting}[language=SQL]
SHUTDOWN;
\end{lstlisting}

\sourcecode{src/apps/minisql/main.ml}{Client command forms.}
\sourcecode{src/apps/minisqld/main.ml}{Server mode selection and startup banner.}

\section{Choose a connection mode deliberately}

\begin{longtable}{p{.32\linewidth}p{.61\linewidth}}
\toprule
Server mode & Meaning and matching client family \\
\midrule
\endhead
\code{--serve} & Trusted loopback. Use \code{--shell}, \code{--query},
\code{--script}, or \code{--ping}. Appropriate for a local tutorial or an
environment whose local access is the intended trust boundary. \\
\code{--serve-authenticated} & Password-authenticated loopback with MiniSQL
protected framing after login. Use \code{--auth-shell-prompt} and the related
prompt commands. This mode is not TLS. \\
\code{--serve-secure} & MiniSQL-specific authenticated encrypted framing.
Use the \code{--secure-*} client family. This is a distinct protocol from TLS. \\
\code{--serve-tls} & Native TLS 1.3 plus database authentication. Use
\code{--tls-*} or \code{--tls-pin-*}. The certificate identity and trust
configuration must match. \\
\code{--serve-standby} & Read-only standby server, bound to loopback. This
requires a database materialized as a standby. \\
\bottomrule
\end{longtable}

Configuration-driven variants and fencing are explained later. The positional
\code{--serve-\allowbreak authenticated} fixes the address to loopback; the configured
authenticated listener has an explicit address from JSON and uses the same
MiniSQL-specific protected framing after authentication. Native TLS remains
the separate \code{--serve-tls} family. Match the client to the actual listener
mode rather than assuming that every encrypted listener speaks TLS.

\section{Organize an installation for repeatable operation}

Separate application binaries, database directories, configuration, keys,
logs, and backup destinations. Give the process write access to the database
and its intended logging/spill locations, read access to required certificates
and encryption keys, and no unnecessary permissions. Use absolute database,
configuration, log, and key paths in unattended launchers. Relative paths
resolve against the process working directory, which can differ between an
interactive shell and a service supervisor.

There is no universal installer wizard in the repository. Copying the native
binaries and supplying an explicit configuration is the deployment model.
The supplied systemd unit and Windows launch/stop scripts are templates to
adapt, not automatic account creation or service registration tools.

Keep a small installation record: source commit, compiler revision, target,
binary directory, database path and logical name, page size, listener mode,
port, configuration path, key-provider paths, and the date of the latest
successful restore drill. This information makes a later upgrade or incident
recovery concrete.

\section{Check the installation before loading important data}

At minimum, verify a query, a committed write, a restart, and an offline check.
Use a disposable database for these checks. A query result proves connection
and execution; a row that survives restart proves a more relevant path for
durable use. A checker result adds storage consistency evidence.

\begin{lstlisting}[language=SQL]
CREATE TABLE installation_check (
  id INTEGER PRIMARY KEY,
  note VARCHAR(80) NOT NULL
);
INSERT INTO installation_check VALUES (1, 'restart verification');
SELECT * FROM installation_check;
SHUTDOWN;
\end{lstlisting}

Run \path{minisql-check <database-path>} while the server is stopped. Start
the server again, query \code{installation\_check}, and confirm the row.
Drop this disposable table after the check if continuing to use the database.

For storage-format qualification, the repository provides a focused gate:

\begin{lstlisting}
.\test.ps1 -StorageV2 -Target windows-x64
.\test.ps1 -StorageV2 -Target linux-x64
\end{lstlisting}

Use the corresponding native host for execution. The direct Python runner
also supports \code{--compiler}, \code{--target}, selected cases,
\code{--integration}, and compile-only cross-target checks. The historical
106-phase release results belong to earlier validated revisions; they do not
automatically certify every new storage-format change. Consult the current
acceptance report and run the relevant gate for the binary being deployed.

\sourcecode{tests/run_storage_v2.py}{Focused storage-v2 compilation and execution gate.}
\sourcecode{docs/acceptance/STORAGE_V2_2026-11-24.md}{Recorded evidence and platform limits.}

\section{Troubleshooting the first connection}

\begin{longtable}{p{.29\linewidth}p{.64\linewidth}}
\toprule
Symptom & Check first \\
\midrule
\endhead
Connection refused & The server must still be running on the selected port.
The shell and server are different programs. \\
Database cannot be opened & Use the actual UUID directory, not only the data
root. Check permissions, the required encryption key, and whether another
process already owns the database lock. \\
Unsupported format & Confirm the source/build revision and persisted format.
Do not edit a header to bypass the check. \\
Authentication failure & Use the matching client mode, an enabled account,
the correct password, and a database \code{CONNECT} grant. \\
TLS handshake failure & Verify TLS 1.3 support, certificate/private-key
availability, the requested server name, and the trust or pin mode. \\
Handshake appears stuck in an older Windows daemon & Exit console selection
mode with Escape. Current code disables QuickEdit before logging, but an
already-running older process may still be paused by text selection. \\
\bottomrule
\end{longtable}

Do not troubleshoot a format or encryption error by deleting metadata, WAL,
lock files, or audit files. Preserve the directory and use the diagnostic
message to identify the failing boundary.
